% Results II — full prose v1.

\subsection{A pre-registered, difficulty-controlled null}

Can the first tokens of a run predict its outcome beyond what the problem
itself predicts? We froze the answer procedure before evaluating it:
endpoints (primary: eventual success of the 16K run; secondary:
scratch-solvability at 4{,}096, i.e.\ $\Rhat(4{,}096) \ge \tauthr$ on
non-instant cells), feature sets (a question-only baseline versus the same baseline plus
fifteen frozen early-window dynamics summaries --- entropy, surprisal,
successive divergences, hidden-state geometry, spectral summaries over the
first 512 tokens; defined in Appendix~\ref{app:signals}), a logistic model, problem-grouped cross-validation on
train+validation, per-class power gates, and the success criterion: the
paired $\Delta$AUROC's problem-clustered 95\% CI must exclude zero. The
question-only baseline contains no signal from the run: the frozen
difficulty columns are the benchmark's human difficulty level, the topic
category, and five surface statistics of the problem text (character
count, whitespace-token count, and counts of numeric, operator, and
equation tokens), plus the model's identity --- everything knowable before
the first token is generated, at text level (a pre-generation
\emph{activation} probe would be a stronger baseline still
\citep{lugoloobi2026}; see Limitations). The test split was then evaluated
once.

Neither endpoint met the criterion (Figure~\ref{fig:conf}). Primary:
baseline AUROC $0.83$ versus $0.85$ with early signals, $\Delta = +0.026$
$[-0.054, +0.167]$. Secondary: baseline $0.88$ versus $0.78$,
$\Delta = -0.090$ $[-0.213, +0.033]$ --- the early features \emph{hurt}
out of sample, the signature of uninformative predictors. A level-free
sensitivity variant agrees, and the within-trajectory timing endpoint
remained below its power gate (four test-split interior events) and was
not fit. A post-hoc sweep of the forecast point over
$t \in \{128, 256, 1{,}024, 2{,}048\}$ --- labeled post-hoc in the
amendment --- found 8 of 10 point estimates negative and every CI
straddling zero (Figure~\ref{fig:sweep}): the null is not an artifact of
the chosen window.

\begin{figure}[t]\centering
\includegraphics[width=\linewidth]{confirmatory_auroc}
\caption{The confirmatory test: single evaluation on the held-out split.
Neither endpoint's $\Delta$AUROC CI excludes zero.}
\label{fig:conf}\end{figure}

\begin{figure}[t]\centering
\includegraphics[width=\linewidth]{early_signal_prefix_sweep}
\caption{Post-hoc forecast-point sweep: no window rescues the early
signals.}
\label{fig:sweep}\end{figure}

Why the null has this shape is visible in the per-model split: the
difficulty baseline alone is near-perfect for Gemma-4 (test AUROC $1.0$
on both endpoints) and mediocre for Ministral-3 ($0.63$--$0.73$).
Difficulty fully determines the capability-limited model's fate, while the
budget-elastic model's outcomes carry residual stochastic variance --- the
band of Section~\ref{sec:results-regimes} --- that neither difficulty nor
early dynamics explain.

\subsection{Why pooled evaluation cannot see within-attempt information}

The confound has a simple formal core. Fix a problem--model cell $c$ and
treat attempts as exchangeable given the cell: attempt $j$ has outcome
$Y_{c,j} \in \{0,1\}$ with $\Pr(Y_{c,j} = 1) = p_c$. Call a predictor
\emph{question-only} if its score depends on the cell alone,
$S_{c,j} = f(c)$; the leave-one-out pass rate
$\hat d_{c,j} = \frac{1}{k_c - 1} \sum_{i \ne j} Y_{c,i}$ is the canonical
example --- unbiased for $p_c$ and independent of attempt $j$'s own
outcome. AUROC is a probability over discordant pairs,
$\mathrm{AUROC}(S) = \Pr(S_a > S_b \mid Y_a = 1, Y_b = 0) + \tfrac{1}{2}
\Pr(S_a = S_b \mid \cdot)$, and splitting the pairs by whether they share
a cell gives
\begin{equation}
\mathrm{AUROC}_{\mathrm{pooled}}(S) \;=\; \lambda \,
\mathrm{AUROC}_{\mathrm{within}}(S) \;+\; (1 - \lambda)\,
\mathrm{AUROC}_{\mathrm{between}}(S),
\label{eq:decomp}
\end{equation}
where $\lambda$ is the fraction of discordant pairs drawn from the same
cell. Two consequences follow. First, every question-only predictor has
$\mathrm{AUROC}_{\mathrm{within}} = \tfrac{1}{2}$ exactly (its score is
constant within a cell, so all within-cell pairs are ties): within-cell
evaluation isolates precisely the information that a difficulty
measurement cannot carry. Second, in the standard design with one attempt
per problem ($k_c = 1$), $\lambda = 0$ identically --- a pooled AUROC then
evaluates only between-cell ranking, a task for which ranking by $p_c$ is
optimal among question-only scores under exchangeability, and is
\emph{structurally incapable} of detecting within-attempt signal, however
high it reads. The published early-window positives are of this form (one trajectory per problem). The two
designs in this section are the two escapes Eq.~\eqref{eq:decomp} allows:
with $k_c = 1$ the within component is inaccessible and the only recourse
is \emph{incremental} value over an explicit question-only baseline (our
confirmatory test above); with $k_c \gg 1$ the within component becomes
directly estimable (the dissection below). One bookkeeping fact used
there: within a cell with $S_c = \sum_i Y_{c,i}$ successes, the LOO score
takes exactly two values, $(S_c - 1)/(k_c - 1)$ for a correct attempt and
$S_c/(k_c - 1)$ for an incorrect one, ranking every incorrect attempt
above every correct one --- its within-AUROC is mechanically $0$, and the
neutral reference for ``no within-attempt information'' is $\tfrac{1}{2}$.

\subsection{The ceiling exists in the literature's own regime}

Our cohort is small, so we asked the same question of public data at
scale. On 192{,}315 DeepSeek-R1 generations over 91{,}573 problems
\citep{openr1math}, a question-only predictor --- the leave-one-out pass
rate of the problem's \emph{other} attempts, for 92\% of problems a single
binary observation --- achieves AUROC $0.873$ $[0.870, 0.876]$
(problem-clustered bootstrap; estimator frozen before evaluation, and
conservative by construction since dataset curation truncates the
difficulty range). That places a question-only signal squarely inside the
$0.79$--$0.95$ range that internal-state probe papers report
\citep{zhang2025knowright, yuan2026diagnostic, sun2026trajectories,
david2025temporal, lugoloobi2026} --- none of which report a question-only
baseline (Figure~\ref{fig:ceiling}). Nor is the ceiling a mathematics
artifact: on 16{,}384 public GPQA-diamond \citep{rein2023gpqa} samples (64 questions $\times$
256 attempts, Llama-3.3-70B \citep{singhi2025solveverify}), the same
estimator reaches $0.917$ $[0.875, 0.938]$ --- despite a four-way
multiple-choice guessing floor that attenuates, never inflates, it.

\begin{figure}[t]\centering
\includegraphics[width=.9\linewidth]{difficulty_ceiling_openr1}
\caption{A question-only signal on public DeepSeek-R1 generations reaches
AUROC 0.873 --- inside the published probe range.}
\label{fig:ceiling}\end{figure}

\subsection{Dissecting a probe positive in the strong-proponents' regime}

Finally, we replicated and then dissected the closest published
early-window positive, using public data and no generation: 128 MATH
problems with 256 R1-Distill-Qwen-7B samples each
\citep{singhi2025solveverify}, re-scored by
teacher-forcing the dump's own tokens through the 4-bit model (top-1
fidelity $0.906$; protocol frozen before extraction). The 256-per-problem
design permits the control no published positive can run: evaluating the
\emph{same} probe within problem, where difficulty is held fixed by
construction.

The uncontrolled result replicates: a last-four-token hidden-state probe
at $t = 4$ reaches pooled AUROC $0.849$ $[0.735, 0.919]$, matching the
$0.84$ reported by \citet{david2025temporal}. The same probe, on the same
trajectories, evaluated within problem --- the $\lambda$-pure component of
Eq.~\eqref{eq:decomp} --- sits at $0.496$ $[0.466, 0.527]$
--- chance --- and remains at chance at every one of ten anchors from
$t = 4$ to $512$ (Figure~\ref{fig:dissect}). The entire positive is
problem identity, far below this dump's 255-sample difficulty ceiling of
$0.974$. Moreover, with sample composition held fixed, the pooled AUROC
itself collapses --- indistinguishable from chance by $t = 16$ ($0.586$,
CI straddling $0.5$) and $\approx 0.55$ from $t = 32$ onward: the
difficulty content readable from last-token states is a prompt echo that
fades as generation proceeds. This reconciles three observations in the literature ---
declining probe signal at larger reasoning budgets \citep{lugoloobi2026},
question-only probes that work before generation
\citep{cencerrado2025noanswer}, and the practice of probing at special
intermediate-\emph{answer} positions rather than arbitrary early positions
\citep{zhang2025knowright} --- under a single reading: hidden states
encode the problem, and increasingly do not encode it, rather than
encoding the fate of the attempt.

\begin{figure}[t]\centering
\includegraphics[width=\linewidth]{math128_probe_dissection}
\caption{Same probe, two evaluations, 6{,}480 public R1-Distill-7B
trajectories: the pooled positive replicates at $t=4$ and is chance within
problem at every anchor.}
\label{fig:dissect}\end{figure}
